% GENERATED FILE. Edit canonical lessons, then run npm run generate:books.
% canonical-source-hash: fnv1a64:b95cdad5a114bff6
% canonical-lessons: GE-C335-diskussion, GE-C335-debatte, GE-C335-vorschlag, GE-C335-vorschlagen, GE-C335-betonen

\chapter{Die Diskussion}
\label{ch:ge-die-diskussion}
\hlchaptermodality{335}

\begin{chapteropening}
\textbf{By the end of this chapter:} \emph{I can take part in a discussion or a debate, make a proposal and stress what matters.}

Complete the last lesson of chapter 335: I can take part in a discussion or a debate, make a proposal and stress what matters.
\end{chapteropening}

\section[die Diskussion]{die Diskussion — the discussion}

\label{lesson:GE-C335-diskussion}



\begin{quote}
Before the new one: say the German for to admit, then the German for to concede.
\end{quote}

\subsection*{You'll want to know: die Diskussion}
\textbf{die Diskussion} — \textquotedblleft{}the discussion\textquotedblright{}.

\begin{quote}
Nach dem Film gab es eine lange Diskussion.
\end{quote}

\begin{quote}
\emph{After the film there was a long discussion.}
\end{quote}

Plural: \emph{die Diskussionen}. It goes with the verb \emph{diskutieren}.

\begin{cousinweb}
\emph{Diskussion} was borrowed from Latin \emph{discussiō} (EtymDB), from \emph{discutere} \textquotedblleft{}to strike apart, to scatter\textquotedblright{}, which later came to mean \textquotedblleft{}to examine\textquotedblright{} (Lewis \& Short).
\end{cousinweb}

\begin{grammarlens}[title={What you've built}]
One more word for telling, explaining and describing at B1.
\end{grammarlens}

\subsection*{Your turn}
\begin{itemize}
\raggedright
  \item \emph{Say it:} \emph{die Diskussion}
  \item \emph{Say it:} \emph{die Diskussion}, once more
  \item \emph{Say it:} a sentence of your own with \emph{die Diskussion}
  \item \emph{Recall:} say the German for to admit, then the German for to concede, then say \emph{die Diskussion} again
\end{itemize}

\begin{tcolorbox}[breakable,colback=teal!4,colframe=teal!35!black,title={Before you move on}]
What does die Diskussion mean? (\textquotedblleft{}The discussion\textquotedblright{}.) Say it once more.
\end{tcolorbox}
\section[die Debatte]{die Debatte — the debate}

\label{lesson:GE-C335-debatte}



\begin{quote}
Before the new one: say the German for to concede, then the German for the discussion.
\end{quote}

\subsection*{You'll want to know: die Debatte}
\textbf{die Debatte} — \textquotedblleft{}the debate\textquotedblright{}.

\begin{quote}
Im Parlament gibt es heute eine Debatte über die Schulen.
\end{quote}

\begin{quote}
\emph{Today there is a debate in parliament about schools.}
\end{quote}

Plural: \emph{die Debatten}. More formal and more public than \emph{die Diskussion}.

\begin{cousinweb}
\emph{Debatte} was borrowed from French; EtymDB links it to \emph{débattre}, Old French \emph{debatre} \textquotedblleft{}to fight, to contend\textquotedblright{}. That verb is built on \emph{battre} \textquotedblleft{}to beat\textquotedblright{}, from Latin \emph{battuere} \textquotedblleft{}to beat\textquotedblright{} (Lewis \& Short).
\end{cousinweb}

\begin{grammarlens}[title={What you've built}]
One more word for telling, explaining and describing at B1.
\end{grammarlens}

\subsection*{Your turn}
\begin{itemize}
\raggedright
  \item \emph{Say it:} \emph{die Debatte}
  \item \emph{Say it:} \emph{die Debatte}, once more
  \item \emph{Say it:} a sentence of your own with \emph{die Debatte}
  \item \emph{Recall:} say the German for to concede, then the German for the discussion, then say \emph{die Debatte} again
\end{itemize}

\begin{tcolorbox}[breakable,colback=teal!4,colframe=teal!35!black,title={Before you move on}]
What does die Debatte mean? (\textquotedblleft{}The debate\textquotedblright{}.) Say it once more.
\end{tcolorbox}
\section[der Vorschlag]{der Vorschlag — the suggestion, the proposal}

\label{lesson:GE-C335-vorschlag}



\begin{quote}
Before the new one: say the German for the discussion, then the German for the debate.
\end{quote}

\subsection*{You'll want to know: der Vorschlag}
\textbf{der Vorschlag} — \textquotedblleft{}the suggestion, the proposal\textquotedblright{}.

\begin{quote}
Ich habe einen Vorschlag: Wir essen heute draußen.
\end{quote}

\begin{quote}
\emph{I have a suggestion: we eat outside today.}
\end{quote}

Plural: \emph{die Vorschläge}. \emph{Einen Vorschlag machen} — to make a suggestion.

\begin{cousinweb}
\emph{Vorschlag} is the noun of \emph{vorschlagen}, built from \emph{vor} \textquotedblleft{}before, forward\textquotedblright{} and \emph{schlagen} \textquotedblleft{}to strike\textquotedblright{}: you put an idea forward. EtymDB has no entry for the compound, so this is the plain reading of its parts.
\end{cousinweb}

\begin{grammarlens}[title={What you've built}]
One more word for telling, explaining and describing at B1.
\end{grammarlens}

\subsection*{Your turn}
\begin{itemize}
\raggedright
  \item \emph{Say it:} \emph{der Vorschlag}
  \item \emph{Say it:} \emph{der Vorschlag}, once more
  \item \emph{Say it:} a sentence of your own with \emph{der Vorschlag}
  \item \emph{Recall:} say the German for the discussion, then the German for the debate, then say \emph{der Vorschlag} again
\end{itemize}

\begin{tcolorbox}[breakable,colback=teal!4,colframe=teal!35!black,title={Before you move on}]
What does der Vorschlag mean? (\textquotedblleft{}The suggestion, the proposal\textquotedblright{}.) Say it once more.
\end{tcolorbox}
\section[vorschlagen]{vorschlagen — to suggest, to propose}

\label{lesson:GE-C335-vorschlagen}



\begin{quote}
Before the new one: say the German for the debate, then the German for the suggestion.
\end{quote}

\subsection*{You'll want to know: vorschlagen}
\textbf{vorschlagen} — \textquotedblleft{}to suggest, to propose\textquotedblright{}.

\begin{quote}
Mein Freund schlägt vor, dass wir zusammen kochen.
\end{quote}

\begin{quote}
\emph{My friend suggests that we cook together.}
\end{quote}

Separable and strong: \emph{er schlägt vor}, \emph{er hat vorgeschlagen}.

\begin{cousinweb}
\emph{vorschlagen} is built from \emph{vor} \textquotedblleft{}forward\textquotedblright{} and \emph{schlagen} \textquotedblleft{}to strike\textquotedblright{}, which is usually traced to Proto-Germanic \emph{slahaną}. The sense of putting an idea forward grew out of the parts; EtymDB has no entry for the compound.
\end{cousinweb}

\begin{grammarlens}[title={What you've built}]
One more word for telling, explaining and describing at B1.
\end{grammarlens}

\subsection*{Your turn}
\begin{itemize}
\raggedright
  \item \emph{Say it:} \emph{vorschlagen}
  \item \emph{Say it:} \emph{vorschlagen}, once more
  \item \emph{Say it:} a sentence of your own with \emph{vorschlagen}
  \item \emph{Recall:} say the German for the debate, then the German for the suggestion, then say \emph{vorschlagen} again
\end{itemize}

\begin{tcolorbox}[breakable,colback=teal!4,colframe=teal!35!black,title={Before you move on}]
What does vorschlagen mean? (\textquotedblleft{}To suggest, to propose\textquotedblright{}.) Say it once more.
\end{tcolorbox}
\section[betonen]{betonen — to stress, to emphasize}

\label{lesson:GE-C335-betonen}



\begin{quote}
Before the new one: say the German for the suggestion, then the German for to suggest.
\end{quote}

\subsection*{You'll want to know: betonen}
\textbf{betonen} — \textquotedblleft{}to stress, to emphasize\textquotedblright{}.

\begin{quote}
Die Lehrerin betont, dass die Prüfung sehr wichtig ist.
\end{quote}

\begin{quote}
\emph{The teacher stresses that the exam is very important.}
\end{quote}

Inseparable: \emph{sie betont, sie hat betont}. First of sounds: \emph{die erste Silbe betonen} — to stress the first syllable.

\begin{cousinweb}
\emph{betonen} is built on \emph{Ton} \textquotedblleft{}tone, accent\textquotedblright{}, which came through Latin \emph{tonus} \textquotedblleft{}a sound, a tone\textquotedblright{} (Lewis \& Short) from Greek \emph{tónos} \textquotedblleft{}a stretching, a tone\textquotedblright{}: you give a word more tone.
\end{cousinweb}

\begin{grammarlens}[title={What you've built}]
One more word for telling, explaining and describing at B1.
\end{grammarlens}

\subsection*{Your turn}
\begin{itemize}
\raggedright
  \item \emph{Say it:} \emph{betonen}
  \item \emph{Say it:} \emph{betonen}, once more
  \item \emph{Say it:} a sentence of your own with \emph{betonen}
  \item \emph{Recall:} say the German for the suggestion, then the German for to suggest, then say \emph{betonen} again
\end{itemize}

\begin{tcolorbox}[breakable,colback=teal!4,colframe=teal!35!black,title={Before you move on}]
What does betonen mean? (\textquotedblleft{}To stress, to emphasize\textquotedblright{}.) Say it once more.
\end{tcolorbox}
